\section{Methodik}
\label{sec:allgemeine_methodik}

Für die Experimente wurde das im Aufbau \ref{sec:aufbau} beschriebene virtuelle Netzwerk genutzt.
Dabei wurde für jedes Experiment eine Datei erzeugt, die den Datenverkehr im Netzwerk aufzeichnet.
iperf3 \ref{sec:iperf3_docker_tool} erzeugt JASON und TXT-Dateien, die direkt die Ergebnisse der Messung enthalten.
Teilweise werden auch Shell Befehle \ref{sec:shell_docker_tool} genutzt, um die Experimente zu starten und zu stoppen oder um Ausgaben der Experimente zu erhalten.
PCAP Dateien werden von dem \ref{sec:pcap_docker_capture} beschriebenen Capture Tool erzeugt und können anschließend mit Wireshark \cite{wireshark} analysiert werden hinsichtlich bestimmter Fragestellungen untersucht werden.
Wireshark ist ein Tool, um Netzwerkverkehr zu analysieren und kann PCAP Dateien öffnen und die enthaltenen Pakete darstellen.
Einzelne Pakete können byteweise analysiert werden.
Wireshark erlaubt einzelne Parameter und Flags der Verbindungen zu analysieren und die Pakete nach bestimmten Kriterien zu filtern.